\chapter{历法}
\chapterkana{暦}
\epigraphbox{“乃命羲和，钦若昊天，历象日月星辰，敬授人时。”}{——《尚书·尧典》}

\kaishou{历法管三件事：下种、动土、杀人。\par
谁定历，谁就定这三件事的日子。\par
后来有一个不识字的人改过一次历。}

\section{十一枚}

城里另有一批文字，管的是外面。这一批也是苏方七十年代初带回来的，与章程总目不同批；编目时合称《规程残卷》，编号 \texttt{ФВ-1971-22}，共存十一枚，枚次不连续，最大的一枚编号到四十四。\footnote{枚次空缺多半是原件缺失，也可能是编号跳号，两种情况在红外成像下分不出来。}

十一枚里，与地面有关的只有四枚。以下按内容排，不按枚次。

第四十四枚最短，只有三行；中间一行是一句设问，那一笔入刀比上下文都深。原刻如下：

\begin{archive}[内部参考]{档案 № ФВ-1971-22}
第四十四枚，共三行。上行：「凡遣物于地，先定其判据。」中行：「土著得历之后，作何状。」下行：「记之，勿断。」
枚末另有一行小字，刀口较浅：「此条自登船前立，未改。」
\end{archive}

「未改」两个字，是这一枚唯一带时间的话。章程总目一千二百四十七条，改动栏全空；这一条不在总目里，也写着未改。两处合起来，能说的只有一句：这套文字从立下到出土，没有被动过。

至于为什么没有被动过，材料里没有理由，本章也不替它编，只把这处空白记在这里，往下走。

\section{可遣之物}

第二枚是一张清单，分两栏：可遣，不可遣。可遣一栏短，四行：历一，青铜，灌溉，医。不可遣一栏长：文字，高阶冶炼，兵器，历之全术。

清单只到目录，细目在第三枚。第三枚逐条注「历一」这个词，注得比别处细：历三百六十条，删三术，余者与城中所用同。三术是岁差、交食、五星逆行。

三百六十条这个数，与后世通行的四分历不合（那是一部一年三百六十五又四分之一日的历），与秦汉初年行用的颛顼历也不合。\footnote{四分历与颛顼历的岁实，见《汉书·律历志》与临沂出土的历书简；两套都不是整三百六十日。}中原的历要分四季、要置闰，一部按三百六十等分排下来的历，在中原年年要调。它合用的地方只有一个：没有四季的地方。

这座城的气候是造的，风、光、温都调过，没有春夏秋冬。\footnote{风、水、气三者如何维持，见《月之都——〈广寒宫残简〉辑佚与考释》论城形制的一章。}城里的人不靠天时下种，也就没有置闰的必要。

若「历之全术」指的是这套历的全部，那么地面拿到的是一份按城里节律排的历；它能定农时，却与天上的实际走得越来越远。这不是疏忽。删掉三术，历照样能用两百年；两百年之后的事，送历的人不在场看。

\section{删去的三术}

三术之中，交食最好验。日食月食可以提前算，算得出，写在历上，到时候对不上，人人都看得见。删掉交食，等于把一件能被当众检验的东西拿掉了。

岁差最难验。它一百年才差一度多，两三代人看不出来，非得有人把几百年的观测记下来，前后一比，才知道冬至点挪了位置。没有文字的地方，这件事做不成。

五星逆行介于两者之间：行星走回头路，肉眼看得出，只是不知其故，删与不删，两三代人都未必说得上来。

三术删尽，历法仍是一部好用的历：定农时够用，祭祀够用，派役够用。它只是不能用来问天上的事。

第三枚的末尾还有一行，写在清单之外：「余一段，无用，留之。」留下的这一段是什么，清单没说。它后来成了本章最后要讲的那件东西，此处只记一笔：删的时候，没有人认出它是什么。

\section{十部与一部}

历法送到中原，约在公元前二千年。第二十七枚是一份观测记录，字小而密，刻痕的槽宽与前面几枚不同。

\begin{archive}[内部参考]{档案 № ФВ-1971-22}
第二十七枚。历既下，夏析之为十，授十方。有穷之君不识字，取其九而焚之，独留其一。枚末一行：「所留者仍其旧，删减未复。」
\end{archive}

这一枚把四件事写实了，历分十部、分授十方、九部被焚、留下的一部仍是删减本；前三件此前只能从神话里推。

十份抄本一字不差，这是正面自己写的。内容既同，合并不增加任何东西：烧掉九部，留下的照样能用，一格不少。所以他要处理的对象在授时权上：九个方国各自定时，才是他要收的东西。一个方国手上有历，就能自定农时、自祀、自纪年；收掉历，这三件事便收到一处。

《淮南子·本经训》把这件事记成「十日并出，羿射九日」；《左传》襄公四年记后羿「因夏民以代夏政」。\footnote{《淮南子》卷一三《本经训》；《左传》襄公四年。两处一为神话，一为史书，记的是同一件事的两面。}十日并出，是十份抄本同时行用；羿射九日，是废掉其中九部。天上多出九个太阳，对农事毫无意义；十个授时中心同时存在，对权力很有意义。神话留下的那个多余的九，正是记录里被焚的九部。

九部的下落，只有「焚之」两个字。有没有人私藏一部，记录里没有。

\section{还原}

第二十七枚的背面另有两行，槽宽比正面粗，笔画也潦草得多，红外成像下才看清：

\begin{archive}[内部参考]{档案 № ФВ-1971-22}
背面，两行。上行：「冬至失期者三，交食不验者一。」下行：「五年而合，其历复全。算法自归，非我授之。」
\end{archive}

第一行记差错，第二行记结果。差错三处冬至、一次交食；结果用了五年。要紧的是「自归」二字：没有人给他，是他自己从观测里算回来的。

补回来的是三术。三术算天，算准了，历就准了；历准了，农时、祭祀、刑戮都跟着准。到这里，他手上有的是一部好历，别的什么都没有。

\begin{figure}[htbp]
\centering
\begin{tikzpicture}[scale=1.0, font=\scriptsize]
  \useasboundingbox (-0.2,-2.0) rectangle (11.4,2.4);
  % 左：删减本
  \node[anchor=west] at (0,2.0) {删减本};
  \draw[line width=0.6pt] (0,1.6) rectangle (3.6,-1.6);
  \node[anchor=west] at (0.2,0.95) {三百六十条};
  \node[anchor=west] at (0.2,0.25) {无岁差};
  \node[anchor=west] at (0.2,-0.45) {无交食};
  \node[anchor=west] at (0.2,-1.15) {无五星逆行};
  % 箭头
  \draw[line width=0.6pt, ->] (3.7,0) -- (5.6,0);
  \node[anchor=south, align=center] at (4.65,0.1) {五年观测\\补回三术};
  % 右：还原本
  \node[anchor=west] at (5.7,2.0) {还原本};
  \draw[line width=0.6pt] (5.7,1.6) rectangle (9.5,-1.6);
  \node[anchor=west] at (5.9,0.95) {三百六十条};
  \node[anchor=west] at (5.9,0.25) {三术俱全};
  \node[anchor=west] at (5.9,-0.45) {与删减前同};
  \node[anchor=west] at (5.9,-1.15) {底层算得出了};
  % 出：密葬法
  \draw[line width=0.6pt, ->] (9.6,0) -- (10.5,0);
  \node[anchor=west] at (10.6,0) {密葬法};
\end{tikzpicture}
\caption{删减本与还原本。删掉的三术补回之后，本来算不出东西的那一层也活了；它不是后加的，原本就在历里；历不全的时候，它算不出东西。}
\figsource{示意。左栏据第三枚，右栏据第二十七枚背面；「密葬法」一名出自月面档案，见本章末。}
\end{figure}

补完之后，他把这部历叫「新历」。在他看来，这就是一部修好了的历。

公元前一九九〇年前后，新历的条文经信道上行，送回月面。档案对这件事只有一句，写在第二十七枚的边栏：「底牌见于人手。」谁的手，没写。

月之都一看就明白：删掉的三术被人补了回来，补得还对得上。一部删减过的历，靠观测一格一格推回了原样。

历法一旦还原，埋在底下的东西就露出来了。那一层从来没有被列进删减的条目：它不像一条算法，倒像几处记数的旧痕，与农时无关，刻的时候便没有动它。它的算法要用岁差、交食、逆行三样做输入：三样缺一样，它算出来的东西就是空的；三样齐了，它才给出结果。

于是「余一段，无用，留之」那八个字有了下文。留下这一段的人不知道它是什么；两百年后，一个不识字的人把它算活了。

这一段能做什么，月面档案有专门的名称：天文密葬法。名字的来历第十章再说，此处只记它的一件本事：天体的位置是可以指定的。

后羿不知道这些。他只知道历修好了，条文写出来了。他手上那部历与删减前的原本已经一样，这件事他自己并不知道。

本章的证据是一批残卷。十一枚，四枚与地面有关，最要紧的两枚是清单与观测记录。残卷没有第二份，也调不出原件之外的东西；凡本章说是推的，都已在行文里标明。

下一章讲两个人：一个写条文，一个不识字而照办。他们各自知道的，与各自不知道的，正好反过来。
